%% kapa: exam=1 round=20 year=2009 period=2 form=A subject=accounting subject_ko=회계학 seq=2 number=2
%% source: Q-Net 감정평가사 기출문제 (cst00309) artlSeq=5024178 file=extracted/1008/5024178/2009년 제20회 감정평가사 1차시험 기출문제/2교시.hwp
%% conversion: kordoc 4.24.51 -> tools/gen1.py (auto); review=proofread-sonnet; flags=none
\begin{question}{2}
다음 중 충당부채 및 우발부채에 대한 회계처리의 내용으로 옳지 \uline{않은} 것은?
\begin{choices}
  \item 충당부채를 발생시킨 사건과 밀접하게 관련된 자산의 처분차익이 예상되는 경우에 당해 처분차익을 차감한 순액으로 충당부채를 인식한다.
  \item 불법적인 환경오염으로 인해 기업이 지불해야 할 환경정화비용은 기업의 미래행위에 관계없이 의무 이행에 자원의 유출이 수반되므로 충당부채를 인식한다.
  \item 어떤 의무에 대하여 제3자와 연대하여 의무를 지는 경우 이행할 의무 중 제3자가 이행할 것으로 기대되는 부분은 우발부채로 처리한다.
  \item 과거사건의 결과로 발표된 경영방침 등을 통하여 기업이 의무를 이행할 것이라는 정당한 기대를 상대방이 가지게 된 때 발생하는 의제의무도 충당부채 인식 대상이 될 수 있다.
  \item 과거에 우발부채로 처리하였더라도 그 이후 상황변화로 인하여 자원의 유출가능성이 매우 높아지고 금액을 신뢰성 있게 추정할 수 있는 경우에는 그러한 가능성의 변화가 발생한 기간에 충당부채로 인식한다.
\end{choices}
\end{question}
\answer{1}
